\chapter{Bridge Geometry and Shear-Stud Solver}
\label{ch:geometry}

\emph{Pull request: \#116.}

\section{Overview}
Up to this point my work had been on the frontend. PR~\#116 moved into the \texttt{core}
backend to compute a piece of real bridge geometry: the \emph{shear studs} that connect the
steel girders to the concrete deck slab in a composite plate-girder bridge. This chapter
describes the code-based geometry formulas I added, the dynamic keys used to store the results,
and a new solver method that acts as the central entry point for derived component geometry.

\section{Motivation}
Previously, shear-stud head dimensions were effectively fixed values. For a design tool that
claims IS/IRC compliance, these dimensions must instead be \emph{derived} from the stud
diameter according to the governing codes, and stored in the output so that both the 3D model
and the design report can display them. The work therefore had three parts: encode the code
formulas, define keys to hold the computed values, and wire a solver that runs the computation
as part of the design.

\section{Geometry Formulas}
The formulas live in a dedicated geometry module,
\texttt{core/bridge\_components/super\_structure/shear\_studs/geometry.py}. Two functions
compute the minimum head diameter and minimum head height of a stud from its shank diameter,
each citing the code clause it implements:

\begin{lstlisting}[caption={Code-based shear-stud head geometry (\texttt{shear\_studs/geometry.py}).},label={lst:studgeom}]
def min_stud_head_diameter(d_stud_mm: float) -> float:
    """Minimum stud head diameter.  [IRC 22:2015 - Cl. 606.6]
       min_head_diameter = 1.5 * d_stud"""
    return round(1.5 * d_stud_mm, 2)

def min_stud_head_height(d_stud_mm: float) -> float:
    """Minimum stud head height.  [IS 3935:1966]
       min_head_height = 0.667 * d_stud"""
    return round(0.667 * d_stud_mm, 2)
\end{lstlisting}

\noindent The two governing relations are therefore
\[
  d_{\text{head,min}} = 1.5\,d_{\text{stud}} \quad\text{[IRC 22:2015, Cl.\ 606.6]},
  \qquad
  h_{\text{head,min}} = 0.667\,d_{\text{stud}} \quad\text{[IS 3935:1966]}.
\]

\section{Dynamic Keys}
So that the computed values can flow through the rest of the application in a consistent,
namespaced way, two keys were added in \texttt{core/utils/common.py}:

\begin{lstlisting}[caption={Namespaced keys for the derived stud geometry (\texttt{core/utils/common.py}).},label={lst:studkeys}]
KEY_DS_STUD_HEAD_DIAMETER = "design_options.shear_studs.head_diameter"
KEY_DS_STUD_HEAD_HEIGHT   = "design_options.shear_studs.head_height"
\end{lstlisting}

Using dotted, namespaced key strings keeps the output dictionary self-describing and makes the
values easy to look up later when building the design report (Chapter~\ref{ch:report}).

\section{The Bridge-Component Solver}
The formulas are invoked from a new solver on the \texttt{PlateGirderBridge} class in
\texttt{core/bridge\_types/plate\_girder/plategirderbridge.py}. I deliberately introduced a
single public entry point, \texttt{bridge\_component\_solver}, that dispatches to per-component
private solvers---starting with \texttt{\_solve\_shear\_studs}---so that future derived
components (bracings, diaphragms, and so on) can be added in the same place without changing the
call site.

\begin{lstlisting}[caption={Central solver dispatch and the shear-stud solver (\texttt{plategirderbridge.py}).},label={lst:solver}]
def bridge_component_solver(self) -> None:
    """Central entry point for derived bridge-component geometry."""
    self._solve_shear_studs()
    # ... placeholders for future derived components ...

def _solve_shear_studs(self) -> None:
    """Compute derived shear-stud geometry and store it in output_dict."""
    d_stud = self.output_dict[KEY_DS_STUD_DIAMETER]
    self.output_dict[KEY_DS_STUD_HEAD_DIAMETER] = min_stud_head_diameter(d_stud)
    self.output_dict[KEY_DS_STUD_HEAD_HEIGHT]   = min_stud_head_height(d_stud)
\end{lstlisting}

\section{Result}
Shear-stud head dimensions are now computed from the stud diameter using the IRC and IS
formulas, stored under clear namespaced keys, and produced by a solver structured for easy
extension. These values feed directly into the 3D geometry and the generated design report.

\osdagfig{figure5.1.png}{Shear studs on the top flange of a girder, with head geometry
derived from the code formulas.}{fig:studs}

\section{Summary}
This chapter added the first backend computation of my internship: code-compliant shear-stud
head geometry, dynamic keys to carry the results, and an extensible
\texttt{bridge\_component\_solver} entry point in the plate-girder module.
